\chapter{Research Backlog}

\section{Phase 1 status: closed}
Phase 1 is scientifically complete. The implemented sequence now covers the full planned baseline/reproduction path:
\begin{enumerate}
    \item P1A--P1B: known-pose IMU--UWB particle-filter baseline and robustness checks;
    \item P1C: unknown-position/yaw global initialization;
    \item P1D: geometry and finite-horizon observability diagnostics;
    \item P1E: Han et al. reproduction audit and explicit implementation contract;
    \item P1F-A--P1F-D: paper-state PF baseline, literal AACOPF mechanism, tuning/freeze, and matched comparison;
    \item P1F-E: rank-deficient geometry negative controls;
    \item P1F-F: non-Gaussian/outlier UWB stress;
    \item P1F-G: scalable bounded-candidate adaptation with explicit diversity guards.
\end{enumerate}
The Phase-1 synthesis is documented in the preceding section and defines the handoff to Phase 2.

\section{Retained Phase-1 decisions}
The following decisions are considered settled for the next phase unless later hardware evidence requires revision.

\paragraph{Known-pose fusion baseline.}
Conventional PF remains the transparent baseline for studying communication-rate and peer-selection effects. The known-pose experiments establish that UWB ranging can correct inertial drift effectively when initialization is sufficiently informative.

\paragraph{Geometry before confidence.}
Auxiliary/peer geometry and unresolved gauge freedoms must be considered before interpreting estimator confidence. Particle concentration, effective sample size, or low range residuals are not sufficient evidence of globally correct localization.

\paragraph{Global acquisition.}
Unknown-pose localization remains a finite-support and mode-survival problem even under informative motion. Larger particle populations increase useful support statistically but do not guarantee monotone or reliable acquisition.

\paragraph{AACOPF interpretation.}
Literal AACOPF is retained only as a reproduced mechanism/reference. It is not treated as a reliable global-acquisition or NLOS-robust solution. P1F-G shows that a bounded and diversity-guarded redistribution can be made scalable and genealogy-safe, and can reduce late error, but it does not improve strict global-pose acquisition over PF.

\paragraph{Measurement robustness.}
Robust UWB treatment is a separate estimator-design axis. If required, it should be introduced explicitly through a robust likelihood, gating, an NLOS/reliability model, or reliability-aware peer/measurement selection rather than through stronger particle copying.

\section{Open questions carried forward}
The following questions remain genuinely open and should not be confused with unfinished Phase-1 work:
\begin{itemize}
    \item \textbf{Real IMU bias/noise:} persistent accelerometer and gyroscope bias can dominate the estimate, but bias-state augmentation should be decided after hardware characterization.
    \item \textbf{Reliable global acquisition:} if broad-prior infrastructure-free global initialization becomes a hard deployment requirement, explicit support creation/rejuvenation may be necessary in addition to safe redistribution.
    \item \textbf{NLOS/outlier handling:} the reliability model and robust measurement mechanism remain to be selected and evaluated.
    \item \textbf{Absolute reference/gauge choice:} infrastructure-free collaboration still requires an explicit reference convention or prior if a globally absolute frame is required.
\end{itemize}

\section{Immediate Phase-2 backlog}
Phase 2 addresses the thesis research question: how much UWB communication and collaboration are required to maintain acceptable localization accuracy? The recommended experiment order is:
\begin{enumerate}
    \item \textbf{P2A -- UWB update-rate sweep:} vary ranging frequency while keeping geometry and peer count controlled. Report localization error, acquisition failures, support/confidence diagnostics, and number of UWB transactions.
    \item \textbf{P2B -- peer-count and geometry sweep:} vary the number of available collaborative nodes and their relative trajectories. Separate the value of additional ranges from the value of improved geometry.
    \item \textbf{P2C -- packet loss and outages:} introduce independent loss, burst outages, and delayed recovery. Measure how quickly uncertainty/error grows and how much communication is required to recover.
    \item \textbf{P2D -- LOS/NLOS and measurement reliability:} introduce controlled NLOS/outlier processes and evaluate an explicit reliability-aware measurement model or gate. Keep this mechanism separate from particle-management adaptations.
    \item \textbf{P2E -- IMU-quality sensitivity:} vary inertial noise/bias using measured hardware values when available; determine whether communication requirements shift with sensor quality.
    \item \textbf{P2F -- communication--accuracy frontier:} combine the Phase-2 sweeps into a common trade-off representation such as position/yaw error versus UWB transactions per second. Once measured radio-energy values exist, extend this to error versus energy.
\end{enumerate}

A fixed-rate/all-peer strategy should be retained as the reference throughout Phase 2. The purpose is not yet to optimize the schedule, but to establish the empirical information/communication frontier that Phase 3 will exploit.

\section{Phase 3 and Phase 4 outlook}
Phase 3 will develop adaptive ranging and peer selection. Candidate policies should use estimator uncertainty, relative geometry/information gain, measurement reliability, and communication/energy cost to decide \emph{when} to range and \emph{which peer} to range. Comparisons must remain against the fixed-rate/all-peer Phase-2 baselines.

Phase 4 will move the selected estimator and communication policy onto the ESP32/DW3000/IMU hardware. It will characterize real UWB and IMU noise, measure energy per ranging transaction, validate timing and embedded computational load, and compare hardware logs with the simulation predictions.

\section{Next concrete milestone}
The next milestone is therefore \textbf{P2A: the controlled UWB update-rate sweep}. No further AACOPF reproduction experiments are planned before that transition.
